% !TeX root = ../papers/nc-p3.tex
\section{被約二十点での五次元印付き空間と六点の割線}
\label{sec:finite-marked-five-space-secants}
\begin{theorem}\label{thm:finite-marked-five-space-secants}
四辺が有限で点スキームが被約二十点、全中央支持4とする。
印 $p$ で二つの異なる四空間 $E_A=w(p_0,-,V_2^*,-)$ と
$E_B=w(V_0^*,-,p_2,-)$ が五空間 $E$ を張れば、
$\mathbf P(E)$ の三行三列の階数一切断 $\Sigma$ は必ず正次元である。
特に反対階数10で両側の和が五次元ずつなら、各側に階数一行列の曲線が必要である。
\end{theorem}
\begin{proof}[十三本以上の割線]
$J=\mathbf P(E_A\cap E_B)$ は平面である。位置2の座標を $A$ により $\mathbf P(E_A)$ と
同定する。平面切断の長さは7以下なので少なくとも13点 $q$ が $J$ の外にある。
印は $J$ に入る。これらの点では $A(q),B(q)$ は独立で、両方とも $p_1,p_3$ から射影した
$q_1,q_3$ を左右核に持つ。よってその割線 $\Lambda_q$ は二行二列行列の空間に入る。
$\Sigma$ が有限と仮定すると行列式がこの直線にちょうど長さ2の部分を切る。
同じ割線は二度使えない。少なくとも $\Lambda_q$ は $\mathbf P(E_B)$ に含まれず、
その交点 $B(q)$ が一点に決まり、位置0の射影は閉埋込みだからである。

有限な $\Sigma$ は $\mathbf P^2\times\mathbf P^2$ の四双一次式の完全交叉で長さ6。
射影四平面を張り、被約なら六つの代表ベクトルの唯一の従属関係には全ての項が現れる。
実際 $(1,1)$ の制限像は五次元の完全跡超平面で、被約因子の跡係数は全て非零である
（定理\ref{thm:finite-marked-secant-pencil}のKoszul計算）。
非被約で全因子が曲線状なら長さ2の部分の台は二点の組または非単純因子の唯一の二重点である。
支持数 $k\le5$ なら総数は $\binom{k}{2}+6-k\le11$。
曲線状でない点があればその接次元は2以下である。
実際Segreの接ベクトルを行と列の二つの変化に分けると、一方の変化が零である
非零接方向は射影四平面内に階数一の直線を作ってしまう。
有限完全交叉なので、接次元2なら二つの局所定義式の位数はともに2以上で長さ4以上となる。
従ってそのような点は一つであり、その点の二重点の割線は全て固定接平面内にある。
この平面内の $A(q)$ は多くとも7点である。相異なる支持間の割線は高々3本、
残る長さ2の因子の接線は高々1本なので13点には達しない。

以下 $\Sigma$ は被約六点とする。もし $k$ 点が $\mathbf P(E_A)$ に入れば $k\le4$ である。
両端が内の割線と両端が外の割線が与える点数はそれぞれ $\binom{k}{2},\binom{6-k}{2}$ 以下。
一端だけが内なら交点 $A(q)$ はその端点自身なので全て合わせて $k$ 点以下となる。
$1\le k\le4$ なら合計は11以下である。よって全六点が外にある。
このとき $J$ の外の各 $A(q)$ に対して $B(q)$ も $J$ の外にあり、
同じ計算で $\mathbf P(E_B)$ も六点を全て避ける。

\end{proof}
\begin{proof}[六点の差の平方関係]
六点のベクトル $z_i$ を $\sum_i z_i=0$ と規格化する。
二超平面の式を $a,b$ とし、$a_i=a(z_i),b_i=b(z_i)$ は全て非零で和は零である。
使用された割線の添字の組を $I$ とする。欠けた組は高々二つである。
$ij\in I$ のとき二つの交点は $z_i/a_i-z_j/a_j$ と $z_i/b_i-z_j/b_j$ により表される。
それぞれの交点の二次Veronese像の線型関係空間を比較する。
定理\ref{thm:finite-square-chart}の反対位置の二次座標像の単元倍での一致により、
二つの関係空間は各組ごとの非零定数倍で同定される。

その関係空間を明示する。六列 $z_i/a_i$ の行列を $U$ とすると核は $k(a_i)$。
差の外積でなく対称平方の関係は $U M U^t=0$、$M\mathbf 1=0$ なる対称行列で表される。
すなわち組 $ij$ の係数は $-M_{ij}$ で、対角成分は行和が零となるように決める。
$M=av^t+va^t$（ここでは $a=(a_i)$ は列）、$\sum v_i=0$ で尽きる。
これは五次元商の対称二乗への射の核が $a$ の対称積全体であるためである。
$v_i=a_i t_i$ とおき、各辺の定数倍を除くと関係空間は
\[
         (t_i+t_j)_{ij\in I},\qquad
         \sum_i a_i t_i=0,\qquad t_i+t_j=0\quad(ij\notin I)
\]
である。$b$ でも同じ表示。全ての辺和を零にすれば $t=0$ なのでパラメータ表示は単射である。
従って $D_0=k^6/\langle e_i+e_j:ij\notin I\rangle$ とすると、
$D_0/k\bar a$ と $D_0/k\bar b$ の間に線型同型があり、
各 $e_i+e_j$ の像を対応する像の非零倍に移す。棒は $D_0$ への像とする。

実は $k\bar a=k\bar b$ である。欠辺がなければ相異なる四添字による
$(e_i+e_j)+(e_k+e_l)=(e_i+e_k)+(e_j+e_l)$ を使う。
外の二添字で $a,b$ は非零なので商でもこの四元の関係は一つだけであり、
同型の四つの倍率は等しい。四添字を動かすと全倍率が等しくなり、
辺和が $D_0$ を張るため主張が従う。
欠辺が $12$ だけ、または $12,13$ なら $e_1=u,e_2=-u$、後者では $e_3=-u$ と置く。
残りの独立基底元を $r,s,\ldots$ と書くと少なくとも三つある。
$(u+r)+(-u+s)=r+s$ の三元は、外の基底元に $\bar a,\bar b$ の非零成分があるので
両商でも唯一の関係を持つ。$r,s$ の全順序対を使えばこれら全ての辺和の倍率が揃う。
それらだけで $D_0$ を張る。
欠辺が $12,34$ なら $D_0$ の基底を $u,v,r,s$ とし $e_1=-e_2=u,e_3=-e_4=v$ とする。
$(\epsilon u+\delta v)+(-\delta v+r)=\epsilon u+r$（$\epsilon,\delta=\pm1$）と
$r$ を $s$ に替えた三元を用いる。こちらも残る $s$ または $r$ の係数が非零で
関係は唯一、全ての符号を動かすと倍率が揃い $D_0$ を張る。
各場合とも、同型がこの生成系の自然な像に共通定数を掛けるので、二つの商の核直線は等しい。

従って欠辺に現れない全添字で $a_i/b_i$ は同じである。そのような添字は少なくとも二つある。
その二つを $i,j$ とすると $ij\in I$ にもかかわらず二つの割線交点は比例し、
$A(q)\notin J$ に反する。これで有限切断を全て除外した。
\end{proof}

\begin{proposition}\label{prop:single-quartic-slice-kernel-symmetry}
四辺が有限で全中央支持4とし、接次元1以下の印において、
一つの四空間 $E_A=w(p_0,-,V_2^*,-)$ の階数一切断が有限長4であるとする。
もう一つの四空間との一致を仮定せず、核側の二位置は反対対称性を持つ。
従って上の被約二十点の五空間で、階数一の整曲線が $\mathbf P(E)$ を張る場合も
核側の反対対称性が成立する。
\end{proposition}
\begin{proof}
$E_A$ だけで長さ4環の積跡による対称表示が得られる。
左右核をその表示で $l,r$ とすれば $A l=A r=0$ が点環全域で成立する。
奇側の非対称を仮定し、定理\ref{thm:finite-artin-kernel-symmetry}の証明で
どの偶座標が必要になるかを確認する。
外積 $l\times r$ の成分は独立で、単元成分を持つ $A$ の列との比例により
局所主イデアル $(\eta)$ を生成する。その支持 $Z$ は $A$ の二次小行列を消すので
$A$ 側の四座標による閉埋込みで長さ4の切断の部分となる。三次持上げに使う指標と
$Z$ 上のベクトルは $A$ の列から決まり、$A$ だけの跡模型で $a l=\alpha 1$ が成立する。
従って同定理の余次元1の三次持上げ、逆数分子の小行列、印での一様化変数による除法は
そのまま適用できる。長さ3および局所 $(1,2,1)$ で使う非零隣接語 $H\chi(a)$、
長さ4の曲線状因子の積での四つの独立な三次式も全て $A$ 側の座標を使う。
同定理のもう一つの偶行列との比較で得ていた反対座標の $Z$ への全射は、
最後の長さ4の共線因子の排除だけに必要であった。現在これは
系\ref{cor:finite-no-four-on-line}が独立に排除する。これで全ての選択肢が矛盾する。

$\mathbf P(E)$ を張る整曲線 $C$ の超平面切断 $C\cap\mathbf P(E_A)$ は
$\mathbf P(E_A)$ を張る。実際切断上消える一次式の $C$ 上の制限は、
切断式で割ると $H^0(\mathcal O_C)=k$ に入る。この切断は従って長さ4以上である。
四空間の階数一軌跡が正次元なら一本の被約母線だけであり、こうした切断は入らない。
有限なら長さ4以下なので等号となり、前半を適用する。
\end{proof}
